\documentclass[11pt]{article}
\usepackage{worksheet_style}

% ---- Toggle: uncomment for filled version ----
%\worksheetfilledtrue
\begin{document}
\worksheetheader{16}{Triple Integrals}

\begin{background}
Everything we did with double integrals extends naturally to three dimensions. Triple integrals compute volume, mass, and moments over solid regions. Cylindrical and spherical coordinates play the same role that polar coordinates did for double integrals---they simplify regions with the right symmetry.
\end{background}

%=============================================================
\wsection{Warm-Up: Double Integral Review}

\textbf{(a)} Evaluate $\ds\int_0^2\int_1^3 xy\,dy\,dx$.

1. Inner ($y$): $\wbml{\left[\frac{xy^2}{2}\right]_1^3 = \frac{9x}{2}-\frac{x}{2} = 4x}$

2. Outer ($x$): $\wbml{\int_0^2 4x\,dx = [2x^2]_0^2 = 8}$

3. Result: $\wbm{\boxed{8}}$

\medskip
\textbf{(b)} Evaluate $\ds\iint_R y\,dA$ where $R$ is bounded by $y = 0$, $x = 0$, and $x+y=1$.

1. Region (Type I): $0 \le x \le 1$, $0 \le y \le 1-x$.

2. Inner: $\wbml{\int_0^{1-x}y\,dy = \frac{(1-x)^2}{2}}$

3. Outer: $\wbml{\int_0^1\frac{(1-x)^2}{2}\,dx = \left[-\frac{(1-x)^3}{6}\right]_0^1 = \frac{1}{6}}$

\medskip
\textbf{(c)} Evaluate $\ds\int_0^{2\pi}\int_0^1 r^2 \cdot r\,dr\,d\theta$.

1. Inner ($r$): $\wbml{\int_0^1 r^3\,dr = 1/4}$

2. Outer ($\theta$): $\wbml{2\pi\cdot\frac{1}{4} = \boxed{\frac{\pi}{2}}}$

%=============================================================
\wsection{Triple Integrals and Fubini's Theorem}

\begin{defbox}{Triple Integral}
\[\iiint_E f(x,y,z)\,dV = \int\!\!\int\!\!\int f(x,y,z)\,dz\,dy\,dx\]
-- When $f = 1$: the triple integral gives the \textbf{volume} of $E$

-- Fubini's theorem: for continuous $f$ on a box $[a,b]\times[c,d]\times[p,q]$, integrate in any order
\end{defbox}

\begin{center}
\tdplotsetmaincoords{60}{110}
\begin{tikzpicture}[tdplot_main_coords,scale=1.2]
    \draw[->] (0,0,0) -- (3,0,0) node[right] {$x$};
    \draw[->] (0,0,0) -- (0,3,0) node[right] {$y$};
    \draw[->] (0,0,0) -- (0,0,3) node[left] {$z$};
    \fill[blue!20,opacity=0.3] (0,0,0) -- (2,0,0) -- (0,2,0) -- cycle;
    \fill[blue!30,opacity=0.3] (0,0,0) -- (2,0,0) -- (0,0,2) -- cycle;
    \fill[blue!30,opacity=0.3] (0,0,0) -- (0,2,0) -- (0,0,2) -- cycle;
    \fill[blue!40,opacity=0.3] (0,0,2) -- (2,0,0) -- (0,2,0) -- cycle;
    \draw[blue,thick] (0,0,0) -- (2,0,0) -- (0,2,0) -- cycle;
    \draw[blue,thick] (0,0,0) -- (0,0,2);
    \draw[blue,thick,dashed] (0,0,2) -- (2,0,0);
    \draw[blue,thick,dashed] (0,0,2) -- (0,2,0);
    \node at (1,3,1.5) {$z = 2-x-y$};
    \node at (0.5,0.5,0) {$z = 0$};
\end{tikzpicture}\\
{\small A typical triple-integral Type I region: the solid tetrahedron in the first octant bounded below by $z=0$ and above by the plane $z = 2-x-y$. Integrate $z$ first from the bottom to top surface, then project onto the $xy$-plane for the remaining double integral.}
\end{center}

\wsubsection{Types of Regions}
\begin{itemize}[nosep, leftmargin=*]
  \item \textbf{Type I:} $z$ bounded by functions of $(x,y)$ -- project onto the $xy$-plane
  \item \textbf{Type II:} $x$ bounded by functions of $(y,z)$ -- project onto the $yz$-plane
  \item \textbf{Type III:} $y$ bounded by functions of $(x,z)$ -- project onto the $xz$-plane
\end{itemize}
Choose the type that makes the bounds simplest.

%=============================================================
\wsection{Examples in Rectangular Coordinates}

\begin{example}[Rectangular Box]
Evaluate $\ds\int_0^1\int_0^2\int_0^3 (x+y+z)\,dz\,dy\,dx$.

\wbbox{
1. \textbf{Inner ($z$):} $\int_0^3(x+y+z)\,dz = \left[(x+y)z+\frac{z^2}{2}\right]_0^3 = 3(x+y)+\frac{9}{2}$

2. \textbf{Middle ($y$):} $\int_0^2\left(3x+3y+\frac{9}{2}\right)dy = \left[3xy+\frac{3y^2}{2}+\frac{9y}{2}\right]_0^2 = 6x+6+9 = 6x+15$

3. \textbf{Outer ($x$):} $\int_0^1(6x+15)\,dx = [3x^2+15x]_0^1 = 18$

4. \textbf{Result:} $\boxed{18}$
}{5cm}
\end{example}

\begin{example}[Type I: First-Octant Tetrahedron]
Find the volume of the tetrahedron bounded by $x = 0$, $y = 0$, $z = 0$, and $x+y+z = 1$.

\wbbox{
1. \textbf{Identify bounds (Type I):}
\begin{itemize}[nosep, leftmargin=*]
  \item $z$: from $0$ to $1-x-y$
  \item $y$: from $0$ to $1-x$
  \item $x$: from $0$ to $1$
\end{itemize}

2. \textbf{Inner ($z$):} $\int_0^{1-x-y}dz = 1-x-y$

3. \textbf{Middle ($y$):} $\int_0^{1-x}(1-x-y)\,dy = \left[(1-x)y-\frac{y^2}{2}\right]_0^{1-x} = \frac{(1-x)^2}{2}$

4. \textbf{Outer ($x$):} $\int_0^1\frac{(1-x)^2}{2}\,dx = \frac{1}{2}\left[-\frac{(1-x)^3}{3}\right]_0^1 = \frac{1}{6}$

5. \textbf{Result:} $\boxed{V = \frac{1}{6}}$

\textit{Check: $V = \frac{1}{6}|a_1(b_2 c_3)| = \frac{1}{3}\cdot\frac{1}{2}\cdot 1\cdot 1\cdot 1 = 1/6$\;\checkmark}
}{5.5cm}
\end{example}

%=============================================================
\wsection{Cylindrical Coordinates}

\begin{formulabox}{Cylindrical Coordinates}
$x = r\cos\theta$, \quad $y = r\sin\theta$, \quad $z = z$
\[dV = r\,dz\,dr\,d\theta\]
-- Use when the region has circular symmetry in $x,y$ (cylinders, cones, paraboloids)

-- Same as polar coordinates in $xy$, with $z$ added
\end{formulabox}

\begin{example}[Cone Volume via Cylindrical]
Find the volume of the solid bounded below by the cone $z = \sqrt{x^2+y^2}$ and above by $z = 3$.

\wbbox{
1. \textbf{Convert:} In cylindrical, the cone is $z = r$. Region: $0 \le r \le 3$, $r \le z \le 3$, $0 \le \theta \le 2\pi$.

2. \textbf{Set up:}
$V = \int_0^{2\pi}\int_0^3\int_r^3 r\,dz\,dr\,d\theta$

3. \textbf{Inner ($z$):} $\int_r^3 r\,dz = r(3-r)$

4. \textbf{Middle ($r$):} $\int_0^3 r(3-r)\,dr = \int_0^3(3r-r^2)\,dr = \left[\frac{3r^2}{2}-\frac{r^3}{3}\right]_0^3 = \frac{27}{2}-9 = \frac{9}{2}$

5. \textbf{Outer ($\theta$):} $2\pi\cdot\frac{9}{2} = 9\pi$

6. \textbf{Result:} $\boxed{V = 9\pi}$

\textit{Check: $V = \frac{1}{3}\pi r^2 h = \frac{1}{3}\pi(9)(3) = 9\pi$\;\checkmark}
}{5.5cm}
\end{example}

%=============================================================
\wsection{Spherical Coordinates}

\begin{formulabox}{Spherical Coordinates}
$x = \rho\sin\phi\cos\theta$, \quad $y = \rho\sin\phi\sin\theta$, \quad $z = \rho\cos\phi$
\[dV = \rho^2\sin\phi\,d\rho\,d\phi\,d\theta\]
-- $\rho$: distance from origin \quad $\phi$: angle from positive $z$-axis ($0 \le \phi \le \pi$) \quad $\theta$: angle in $xy$-plane

-- Use when the region is a sphere/hemisphere or integrand involves $x^2+y^2+z^2$

-- Key relation: $x^2+y^2+z^2 = \rho^2$
\end{formulabox}

\begin{example}[Integral over a Sphere]
Evaluate $\ds\iiint_E (x^2+y^2+z^2)\,dV$ where $E$ is the ball $x^2+y^2+z^2 \le 4$.

\wbbox{
1. \textbf{Convert:} $x^2+y^2+z^2 = \rho^2$. Bounds: $0 \le \rho \le 2$, $0 \le \phi \le \pi$, $0 \le \theta \le 2\pi$.

2. \textbf{Set up:}
$\int_0^{2\pi}\int_0^\pi\int_0^2 \rho^2\cdot\rho^2\sin\phi\,d\rho\,d\phi\,d\theta$

3. \textbf{Separate (all bounds are constants):}
$= \left(\int_0^{2\pi}d\theta\right)\left(\int_0^\pi\sin\phi\,d\phi\right)\left(\int_0^2\rho^4\,d\rho\right)$

4. \textbf{Evaluate each:} $(2\pi)(2)\left(\frac{32}{5}\right) = \frac{128\pi}{5}$

5. \textbf{Result:} $\boxed{\dfrac{128\pi}{5}}$
}{5cm}
\end{example}

%=============================================================

\begin{practice}

\prob Evaluate $\ds\int_0^1\int_0^1\int_0^1 (x+y+z)\,dz\,dy\,dx$.

\wans{\wbm{$3/2$}}

\vspace{3.5cm}

\prob Evaluate $\ds\iiint_E z\,dV$ where $E$ is bounded below by $z = 0$ and above by the paraboloid $z = 4-x^2-y^2$. (Use cylindrical.)

\wans{\wbml{$32\pi/3$}}

\vspace{5cm}

\prob Find the volume of the upper hemisphere $x^2+y^2+z^2 \le 9$, $z \ge 0$ using spherical coordinates.

\wans{\wbm{$18\pi$}}

\vspace{4.5cm}

\prob Evaluate $\ds\iiint_E \sqrt{x^2+y^2+z^2}\,e^{-(x^2+y^2+z^2)}\,dV$ over all of $\R^3$. (Use spherical.)

\wans{\wbm{$2\pi$}}

\vspace{5cm}

\end{practice}

\end{document}
